%!TEX root = ../main.tex
\chapter{Conclusiones}
\label{ch:conclusiones}

En este capítulo se presentan las conclusiones del trabajo realizado, se 
discuten las contribuciones y limitaciones del diseño propuesto, y se delinean 
posibles extensiones para trabajo futuro.

\section{Resumen de resultados}

En esta tesis se diseñó un sistema de control para la regulación de la 
velocidad angular del rotor en un aerogenerador basado en un generador 
síncrono de imanes permanentes (PMSG), operando en la región 2 de operación. 
El controlador propuesto combina una estructura de controladores PI en cascada 
con compensación \textit{feedforward} de los términos de acoplamiento cruzado 
entre los ejes $d$ y $q$ del generador, junto con un algoritmo de seguimiento 
del punto de máxima potencia (MPPT) basado en la estrategia de perturbación y 
observación (P\&O).
\\

El desarrollo del trabajo siguió una metodología sistemática que partió del 
modelado matemático del sistema, pasando por el diseño del controlador y el 
análisis formal de estabilidad, hasta la validación mediante simulación 
numérica. Los principales resultados obtenidos se resumen a continuación.
\\

Se obtuvo el modelo dinámico completo del sistema de conversión de energía 
eólica, acoplando el subsistema mecánico del rotor con el subsistema eléctrico 
del PMSG en coordenadas $dq$. A partir de este modelo se formularon los 
objetivos de control: mantener $i_d = 0$ para minimizar pérdidas eléctricas y 
regular $\omega_m$ hacia la velocidad angular óptima determinada por el 
algoritmo MPPT.
\\

Se demostró que la sustitución de las leyes de control en las ecuaciones 
dinámicas del sistema produce un desacoplamiento del sistema en lazo cerrado 
en dos subsistemas independientes, $\Sigma_1$ y $\Sigma_2$ (Lema \ref{lem:wecs}). Este 
resultado es consecuencia directa de la cancelación \textit{feedforward} de los 
términos no lineales $pL i_q \omega_m$ y $-pL i_d \omega_m + p\varphi\omega_m$ 
en las leyes de control $v_d$ y $v_q$, respectivamente. El desacoplamiento 
permitió analizar la estabilidad de cada subsistema de manera independiente.
\\

Para el subsistema lineal $\Sigma_1$, se obtuvieron condiciones necesarias y 
suficientes de estabilidad mediante el criterio de Routh-Hurwitz aplicado a su 
polinomio característico de segundo orden: $k_{P1} > -R$ y $k_{I1} > 0$ 
(Lema \ref{lem:sigma_1}). Estas condiciones garantizan estabilidad exponencial global del 
punto de equilibrio $\bar{x} = 0$, lo cual asegura que la corriente $i_d$ 
converge a cero independientemente de las condiciones iniciales.
\\

Para el subsistema no lineal $\Sigma_2$, se procedió a linealizar alrededor del 
punto de equilibrio deseado $\bar{z}$ (Lema \ref{lem:sigma_2_eq}) y se aplicó el criterio de 
Routh-Hurwitz al polinomio característico de cuarto orden resultante. Se 
obtuvieron cuatro condiciones sobre las ganancias $k_{P2}$, $k_{I2}$, $k_{P3}$ 
y $k_{I3}$ (Lema \ref{lem:sigma_2}), y se propuso un algoritmo sistemático para la 
selección de ganancias que satisfagan dichas condiciones (Lema \ref{lem:conds}). El 
análisis del lugar geométrico de las raíces complementó este resultado, 
visualizando la migración de los valores propios de $\bar{A}_2$ desde el 
semiplano derecho hacia el semiplano izquierdo conforme $k'_{P2}$ aumenta, lo 
cual valida geométricamente la existencia de un umbral mínimo de ganancia para 
la estabilidad.
\\

Finalmente, las simulaciones en MATLAB/Simulink confirmaron el correcto 
funcionamiento del controlador. Ante una entrada escalón de viento constante 
($v_w = 11$ m/s), el sistema alcanzó el estado estable en aproximadamente 
30 s con un sobrepico de $\approx 6$ rad/s en la velocidad angular. Ante un 
perfil de viento turbulento con velocidades entre 9.8 y 16 m/s, el controlador 
demostró capacidad de seguimiento de la referencia variable $\omega_m^\star$ 
generada por el MPPT, manteniendo la corriente $i_d$ regulada en cero durante 
toda la simulación. La eficiencia media de extracción de potencia obtenida fue 
de $\eta_{\text{avg}} = 99.8$ \%, lo cual indica que el sistema opera de 
manera consistente cerca del punto de máxima potencia.

\section{Contribuciones y limitaciones}

El controlador propuesto en esta tesis se sitúa metodológicamente entre dos 
antecedentes directos: el controlador PI clásico de Chinchilla et 
al. \cite{chinchilla2006} y el controlador PI+PBC de Cisneros et 
al. \cite{cisneros2020}.
\\

El esquema de Chinchilla et al.\ emplea la misma arquitectura de lazos PI en 
cascada con compensación \textit{feedforward} de los acoplamientos $dq$ que se 
adopta en la presente tesis. Sin embargo, el diseño de Chinchilla et al.\ se 
basa exclusivamente en técnicas de dominio frecuencial aplicadas a cada lazo de 
corriente de manera aislada, modelando la planta como una función de 
transferencia de primer orden $1/(R_s + L_s s)$ y sin considerar la dinámica 
mecánica del rotor acoplada al generador. En consecuencia, dicho trabajo no 
ofrece garantías formales de estabilidad para el sistema completo en lazo 
cerrado.
\\

Por otra parte, Cisneros et al.\ proponen un controlador PI+SPBC que explota 
las propiedades de pasividad del sistema electromecánico modelado en forma 
Euler-Lagrange para demostrar estabilidad asintótica del sistema completo. La 
prueba de estabilidad se construye mediante una función de Lyapunov de tipo 
energético y el teorema de LaSalle-Yoshizawa, lo que proporciona garantías de 
estabilidad global. No obstante, la implementación del PI+SPBC es 
considerablemente más compleja: requiere la medición o estimación del torque 
mecánico $T_m$ y su derivada $T_m'(\omega_m)$, el cálculo en línea de señales 
auxiliares como $x_{d2}$ y $\varphi$, y la integración de una dinámica 
adicional ($x_{d4}$) para manejar la incertidumbre en el voltaje del bus DC. 
Además, el sistema considerado en \cite{cisneros2020} incluye el rectificador 
y la carga del bus DC, lo que amplía el alcance del modelo pero también la 
complejidad del controlador.
\\

La contribución principal de la presente tesis consiste en proporcionar un 
análisis riguroso de estabilidad para la estructura PI clásica de Chinchilla 
et al., considerando el sistema acoplado completo (mecánico + eléctrico) en 
lazo cerrado. En particular:

\begin{enumerate}
    \item Se demostró formalmente el desacoplamiento del sistema en lazo cerrado 
    en los subsistemas $\Sigma_1$ y $\Sigma_2$, resultado que 
    en \cite{chinchilla2006} se asume implícitamente pero no se demuestra.

    \item Se derivaron condiciones necesarias y suficientes sobre las ganancias 
    del controlador que garantizan estabilidad exponencial del punto de 
    equilibrio deseado, incluyendo el efecto de la no linealidad del torque 
    mecánico $T_m(\omega_m)$ a través de la constante $\gamma$ en la 
    linealización del subsistema $\Sigma_2$.

    \item Se propuso un algoritmo sistemático para la selección de ganancias 
    que satisfagan las condiciones de Routh-Hurwitz, junto con una demostración 
    de existencia de solución (Lema \ref{lem:conds}).

    \item Se complementó el análisis con el lugar geométrico de las raíces, 
    proporcionando una herramienta visual y complementaria a las condiciones 
    algebraicas para la sintonización de ganancias.
\end{enumerate}

Este enfoque ofrece un compromiso entre la simplicidad de implementación del 
controlador PI clásico y el rigor teórico del enfoque basado en pasividad, 
utilizando herramientas de análisis de control lineal (Routh-Hurwitz, lugar 
geométrico de las raíces) que resultan más accesibles para la práctica de la 
ingeniería de potencia.
\\

No obstante, el trabajo presenta limitaciones que deben señalarse. La principal 
es que el análisis de estabilidad se basa en la linealización del sistema 
alrededor del punto de equilibrio (método indirecto de Lyapunov), lo cual 
proporciona garantías de estabilidad \textit{local}. En contraste, el enfoque 
de Cisneros et al.\ basado en pasividad ofrece estabilidad \textit{global}. 
Esto significa que las garantías de estabilidad derivadas en esta tesis son 
válidas en una vecindad del punto de equilibrio cuyo tamaño no se cuantifica 
explícitamente.
\\

Adicionalmente, la suposición de velocidad del viento constante ($v_w$ 
constante) es necesaria para el análisis de estabilidad, dado que el sistema 
$\Sigma_2$ resulta invariante en el tiempo únicamente bajo esta hipótesis. 
Aunque las simulaciones demuestran un desempeño adecuado ante perfiles de 
viento turbulento, no se dispone de una prueba formal de estabilidad para el 
caso de viento variable.
\\

Finalmente, no se evaluó la robustez del controlador ante variaciones 
paramétricas. En \cite{cisneros2020} se demostró mediante simulación que el 
controlador PI clásico de \cite{chinchilla2006} pierde la capacidad de regular 
la velocidad angular ante variaciones simultáneas del 300\,\% en el flujo 
magnético $\varphi$, del 70\,\% en la inductancia $L$ y del 95\,\% en el 
voltaje del bus DC, mientras que el PI+SPBC mantiene su desempeño. Dado que el 
controlador propuesto en esta tesis comparte la misma estructura del PI clásico, 
es razonable esperar una sensibilidad similar ante variaciones paramétricas de 
gran magnitud, lo cual constituye una limitación inherente al enfoque basado en 
linealización y cancelación de términos no lineales.

\section{Trabajo futuro}

A partir de las limitaciones identificadas y de los resultados obtenidos, se 
proponen las siguientes líneas de trabajo futuro:
\\

\begin{enumerate}
    \item \textbf{Extensión a la región 3 de operación.} El controlador diseñado 
    en esta tesis se limita a la región 2, donde el objetivo es maximizar la 
    extracción de potencia. En la región 3, donde la velocidad del viento supera 
    el valor nominal y se requiere limitar la potencia para proteger los 
    componentes de la turbina, es necesario incorporar estrategias adicionales 
    como el control de ángulo de paso (\textit{pitch control}). Una extensión 
    natural del presente trabajo consistiría en diseñar un controlador que 
    integre el control de torque con el control de \textit{pitch} para cubrir 
    el rango completo de operación del aerogenerador, incluyendo las 
    transiciones entre regiones.

    \item \textbf{Validación experimental.} Todos los resultados presentados en 
    esta tesis se obtuvieron mediante simulación numérica. Una validación 
    experimental del controlador en un prototipo de laboratorio ---similar a la 
    reportada en \cite{chinchilla2006} con un motor DC que emula la turbina 
    eólica--- permitiría verificar el desempeño del controlador ante condiciones 
    reales de operación, incluyendo ruido en las mediciones, retardos de 
    cómputo, efectos de la modulación PWM y saturación de los actuadores.
\end{enumerate}
